\documentclass[12pt]{article}
\usepackage[utf8]{inputenc}
\usepackage[T1]{fontenc}
\usepackage{hyperref}
\title{Tobati Language Resources for K-12 Curriculum Development}
\author{Compiled by Claude Sonnet 5 (claude-sonnet-5) for Andr\'as Kornai}
\date{}
\begin{document}
\maketitle
\section{Introduction}
Tobati (ISO 639-3 tti; also Yotafa) is an Oceanic Austronesian language of the Sarmi--Jayapura subgroup, spoken by the Tobati people around Youtefa Bay, Jayapura, Papua. It is severely endangered; sources give very different speaker counts (about 100 in 2007; another report says only six fluent speakers, all over 60). No literacy is reported.
\section{Primary Educational Resources}
None found in search results.
\section{Linguistic and Reference Materials}
\begin{itemize}
\item Wikipedia, Tobati language: \url{https://en.wikipedia.org/wiki/Tobati_language}
\item Wikipedia, Tobati people: \url{https://en.wikipedia.org/wiki/Tobati_people}
\item Numeral system survey (MPI-EVA): \url{https://lingweb.eva.mpg.de/channumerals/Tobati.htm}
\item Endangered Languages Project entry: \url{https://endangeredlanguages.com/lang/tti?hl=en}
\end{itemize}
\section{Religious and Community Resources}
\begin{itemize}
\item Papua Heritage article on disappearing Papuan languages: \url{https://papuaheritage.org/id/node/35}
\item Good News from Indonesia, article on Kampung Tobati: \url{https://www.goodnewsfromindonesia.id/2022/02/12/kampung-tobati-desa-wisata-di-jayapura-dengan-keunikan-hutan-perempuan}
\end{itemize}
\section{Digital and Audio Resources}
A radio programme page (VOI/RRI) on the Tobati language was listed but not opened.
\begin{itemize}
\item \url{https://voinews.id/index.php/component/k2/item/485-tobati-language}
\end{itemize}
\section{Curriculum Development Context}
Tobati is not known to be taught in schools, and with so few elderly speakers, documentation must come before any curriculum. No orthography, dictionary or teaching material was found.
\end{document}